% =====================================================================
% Глава 2. Постановка задачи и методология исследования.
% Объем 5--6 страниц.
% =====================================================================

\section{Постановка задачи и методология исследования}
\label{sec:methodology}

Глава формализует объект исследования, исследовательские вопросы
и методологию их эмпирической проверки --- математическая модель,
шесть топологий, пять ролей человека, система метрик, обоснование
корпуса и моделей, воронка из четырех экспериментов.

\subsection{Формальная модель мульти-агентной системы}
\label{sec:meth:formal}

Мульти-агентная система рассматривается как кортеж
\begin{equation}
\mathcal{M} = \langle V,\, E_t,\, R,\, \mathcal{T},\, F,\, h \rangle,
\end{equation}
где $V = \{v_1, \dots, v_n\}$ --- конечное множество агентов,
каждый из которых реализован как языковая модель с фиксированной
системной подсказкой и набором инструментов; $E_t \in 2^{V \times V}$
--- множество допустимых направленных ребер коммуникации на
итерации~$t$, образующее коммуникационный граф; $R \colon V \to
\mathcal{R}_A$ --- отображение агентов в множество ролей агентов
$\mathcal{R}_A$ из шести значений (планировщик, исследователь,
исполнитель, критик, дебатер, координатор); $\mathcal{T}$ ---
текущая задача с целевой функцией оценки качества решения;
$F$ --- конечное множество из четырех фаз решения задачи
(планирование, исполнение, проверка, завершение); $h$ ---
(опционально) участник-человек с~ролью из~отдельного множества
$\mathcal{R}_H$ (см.~\ref{sec:meth:roles}).

В классических работах граф $E_t$ определяется архитектурно и не
меняется в ходе выполнения задачи. В настоящем исследовании
вводится расширение модели --- структура $E_t$ становится результатом
оператора обновления, зависящего от~текущей фазы $f_t \in F$ и
наблюдаемых сигналов $\sigma_t \in \{0,1\}^m$ агентов:
\begin{equation}
E_t = \Phi(f_t,\, \sigma_t,\, t),
\label{eq:meth:adaptive}
\end{equation}
где $\Phi$ --- оператор обновления структуры коммуникационного
графа, $m$ --- фиксированная размерность вектора сигналов
(флаги завершения этапа, признаки тупикового состояния, счетчики
отказов критика и другие наблюдаемые индикаторы прогресса).
Возможность такого переключения и составляет суть исследуемой
адаптивной мета-топологии.

Множество ролей агентов $\mathcal{R}_A$ выше уже зафиксировано
шестью значениями. Множество ролей человека~$\mathcal{R}_H$
фиксировано и содержит пять значений, формализованных в
разделе~\ref{sec:meth:roles}. Фазы $F$ упорядочены отношением
$\textit{planning} \prec \textit{execution} \prec
\textit{verification} \prec \textit{done}$ с монотонной семантикой
переходов.

Задача оптимизации мульти-агентной системы формулируется как
многокритериальная по~двум первичным осям~--- качеству решения
$q(\mathcal{M})$ и~суммарной стоимости вызовов $c(\mathcal{M})$~---
при~ограничении на~время выполнения
$\tau(\mathcal{M}) \le \tau_{\max}$, где $\tau_{\max}$ задается
конфигурацией эксперимента. Поскольку чистый Парето-оптимум обычно
не единственен, в работе используется скаляризация с фиксированными
порогами улучшения по~двум первичным осям. Конфигурация
$\mathcal{M}_1$ считается выигрышной относительно $\mathcal{M}_0$,
если соблюдено ограничение $\tau(\mathcal{M}_1) \le \tau_{\max}$ и
выполняется одно из условий
\begin{equation}
\frac{q(\mathcal{M}_1) - q(\mathcal{M}_0)}{q(\mathcal{M}_0)}
   \geq 0{,}05
   \quad\text{или}\quad
\frac{c(\mathcal{M}_0) - c(\mathcal{M}_1)}{c(\mathcal{M}_0)}
   \geq 0{,}15
   \text{ при } q(\mathcal{M}_1) \geq q(\mathcal{M}_0).
\label{eq:meth:thresholds}
\end{equation}
Пороги $5\%$ для качества и $15\%$ для стоимости зафиксированы
до проведения основных запусков и обоснованы в
разделе~\ref{sec:meth:metrics}.

\subsection{Шесть рассматриваемых топологий}
\label{sec:meth:topologies}

Пространство сравнения составляют шесть топологий. Пять из них
являются статичными и реализуют пять базовых классов коммуникационных
структур, перечисленных в обзоре литературы. Шестая является
мета-топологией, переключающей активную статическую топологию
в зависимости от фазы.

\textbf{Star} --- центральный координатор $v_c$ последовательно
обращается к специализированным работникам $v_1, \dots, v_k$ и
агрегирует их ответы. Граф ребер $E_{\text{star}} =
\{(v_c, v_i), (v_i, v_c) \mid i = 1, \dots, k\}$.

\textbf{Chain} --- линейный конвейер из трех агентов
$v_p \to v_e \to v_c$ (планировщик, исполнитель, критик)
с условным циклом доработки. Граф направленный ациклический
с обратной связью от критика к исполнителю.

\textbf{Mesh} --- полносвязная топология с общей шиной сообщений.
Граф $E_{\text{mesh}} = (V \times V) \setminus \{(v,v)\}$, обмен
организован по принципу циклической очереди, решение принимается
голосованием.

\textbf{Debate} --- пара оппонентов $v_+$ и $v_-$ генерирует
параллельные альтернативные решения, после чего судья $v_j$
выбирает победителя или синтезирует ответ. Граф двудольный
с агрегатором.

\textbf{Hierarchical} --- двухуровневая иерархия из координатора
верхнего уровня $v_c^{\text{top}}$ и двух подкоманд
$T_a, T_b \subset V$, каждая из которых сама по себе является
графом. Подкоманды работают параллельно, координатор синтезирует
итог.

\textbf{Adaptive} --- мета-топология, активирующая на каждой
итерации одну из пяти базовых конфигураций. Решение о выборе
принимается маршрутизатором согласно
формуле~(\ref{eq:meth:adaptive}). Псевдокод процедуры приведен
в таблице~\ref{tab:meth:adaptive}, общая концептуальная схема
взаимодействия блоков системы изображена на
рисунке~\ref{fig:arch:functional} главы~\ref{sec:architecture},
архитектура реализации описана в разделе~\ref{sec:arch:topology}.

\begin{table}[!htbp]
\caption{Процедура работы адаптивной мета-топологии}
\label{tab:meth:adaptive}
\small
\begin{tabularx}{\textwidth}{r X}
\toprule
\textbf{Шаг} & \textbf{Действие} \\
\midrule
1  & Инициализировать фазу $f \leftarrow planning$, начальную
     топологию $K \leftarrow K_0$, состояние $s \leftarrow s_0$
     и журнал $\mathcal{J} \leftarrow \varnothing$. \\
2  & Пока $f \neq done$ и $c(s) < B$ выполнять шаги 3--10. \\
3  & Применить маршрутизатор фаз
     $f' \leftarrow \mathrm{PhaseRouter}(s, f)$. \\
4  & Применить маршрутизатор топологий
     $K' \leftarrow \mathrm{TopologyRouter}(s, K, f')$. \\
5  & Если $\mathrm{Guard}(s, K, K') = \mathit{block}$, то
     $K' \leftarrow K$ и записать $\mathit{guard\_override}$
     в $\mathcal{J}$. \\
6  & Если $K' \neq K$, применить ворота перехода состояния
     $s \leftarrow \mathrm{TransitionGate}(s, K, K')$ и записать
     переход в $\mathcal{J}$. \\
7  & Исполнить подграф выбранной топологии
     $s \leftarrow \mathrm{Subgraph}_{K'}(s)$. \\
8  & Обновить $K \leftarrow K'$ и $f \leftarrow f'$. \\
9  & Записать журнал итерации в $\mathcal{J}$. \\
10 & Перейти к шагу 2. \\
11 & Вернуть результат $y(s)$ и журнал $\mathcal{J}$. \\
\bottomrule
\end{tabularx}
\end{table}

\subsection{Пять ролей человека в контуре управления}
\label{sec:meth:roles}

Пять ролей человека формализованы по аналогии с типовыми позициями
в человеческих организациях. \textbf{Координатор} ($\rho_C$) ---
управляет планом и делегированием, активен в иерархических и
полносвязных топологиях. \textbf{Рецензент} ($\rho_R$) --- проверяет
промежуточные результаты и инициирует доработку, естественно
вписывается в Star и Chain. \textbf{Судья} ($\rho_J$) ---
выносит финальный вердикт при наличии альтернатив, релевантен
для Debate. \textbf{Равноправный участник} ($\rho_P$) ---
вносит альтернативные мнения наравне с агентами, вписывается
в Mesh. \textbf{Наблюдатель} ($\rho_M$) --- следит за процессом
и вмешивается только в критических случаях (превышение бюджета,
зависание, циклы), применим в любой топологии.

Не все пары (топология, роль) являются осмысленными.
Рассматриваются пять статических топологий и пять ролей человека ---
полное произведение содержит 25 пар. Из него исключаются три
заведомо вырожденные комбинации. Пара (Debate, Coordinator)
исключается, поскольку координатор в дебатной топологии
эквивалентен судье. Пары (Chain, Peer) и (Hierarchical, Peer)
исключаются, поскольку равноправный участник ломает линейную или
иерархическую структуру. После исключений пространство содержит
$5 \times 5 - 3 = 22$ осмысленные комбинации
$\text{топология} \times \text{роль человека}$, где~используются
пять статических топологий и~пять ролей человека из~$\mathcal{R}_H$.

\subsection{Система оценочных метрик}
\label{sec:meth:metrics}

Метрики разделены на три группы. Первая группа --- базовые метрики
качества и эффективности --- применима ко всем экспериментам и
ко всем топологиям.

\begin{itemize}[noitemsep]
  \item $q \in [0, 1]$ --- агрегированный показатель качества
        решения. Конкретная формула зависит от класса задач ---
        для задач программирования это метрика pass@1 (доля
        задач, для которых сгенерированный код проходит все
        предзаданные модульные тесты с первой попытки); для
        математических задач --- exact match (точное совпадение
        числового ответа с эталоном); для творческой генерации ---
        полусумма метрики ROUGE-L~\cite{rouge2004} (от англ.
        Recall-Oriented Understudy for Gisting Evaluation, Longest
        common subsequence variant) и~доли покрытых концептов; для
        анализа данных --- exact match числового ответа.
  \item $c$ --- суммарная стоимость вызовов языковых моделей
        в долларах США за один запуск.
  \item $\tau$ --- время выполнения запуска по настенным часам
        (от англ. wall-clock) в секундах.
  \item $r_{\text{fail}}$ --- доля запусков, завершившихся со
        статусом, отличным от успешно выполненного.
  \item $N_{\text{iter}}$ --- число итераций основного цикла.
\end{itemize}

Вторая группа --- метрики, специфичные для адаптивной мета-топологии.

\begin{itemize}[noitemsep]
  \item $N_{\text{switch}}$ --- число фактических переключений
        активной топологии за один запуск.
  \item $r_{\text{guard}}$ --- доля решений маршрутизатора,
        заблокированных защитными правилами.
  \item $\gamma$ --- доля бюджета, потраченная маршрутизатором
        на собственные вызовы языковой модели; критичная метрика,
        отвечающая на вопрос, не превышает ли стоимость адаптации
        получаемый от нее выигрыш.
  \item $r_{\text{hurt}}$ --- доля задач, на которых адаптивная
        конфигурация дает качество хуже лучшей статической
        конфигурации. Является ключевым индикатором безопасности
        механизма адаптации.
  \item $\Delta_{\text{oracle}}$ --- разрыв качества между
        теоретическим per-task best-of референсом и~фактическим
        маршрутизатором, формально определяемый
        по~формуле~(\ref{eq:meth:delta-oracle}). Per-task best-of
        референс вычисляется по~результатам базового эксперимента~E1~---
        для~каждого корпуса задач выбирается статическая топология
        с~максимальным средним качеством на~этом корпусе. Референс
        служит верхней границей достижимого качества для~класса
        задач при~заданном пространстве статических топологий и~требует
        апостериорного знания пары $(\text{корпус},\text{лучшая топология})$;
        в~условиях открытого развертывания без~классификатора задач недостижим.
\end{itemize}

\begin{equation}
\Delta_{\text{oracle}} = q_{\text{best-of}} - q_{\text{router}},
\quad\text{где}\quad
q_{\text{best-of}} = \frac{1}{|\mathcal{C}|}
   \sum_{c \in \mathcal{C}}
      q\bigl(K^*_c,\, c \bigr),
\label{eq:meth:delta-oracle}
\end{equation}

\noindent где $\mathcal{C}$ --- множество корпусов задач,
$K^*_c = \arg\max_K q(K, c)$ --- оптимальная статическая топология
для~корпуса~$c$, $q(K, c)$ --- среднее качество топологии~$K$
на~корпусе~$c$ по~базовому эксперименту~E1.

Третья группа --- метрики, относящиеся к взаимодействию с человеком.

\begin{itemize}[noitemsep]
  \item $N_{\text{int}}$ --- число обращений системы к человеку
        за один запуск.
  \item $L_{\text{cog}}$ --- прокси-метрика когнитивной нагрузки
        для запусков с симулятором человека, формально определяемая
        по формуле~(\ref{eq:meth:lcog}). В экспериментах с реальными
        участниками заменяется на стандартную шкалу NASA-TLX.
\end{itemize}

\begin{equation}
L_{\text{cog}} = w_n \cdot \widetilde{N}_{\text{int}}
              + w_l \cdot \widetilde{L}_{\text{ctx}}
              + w_\tau \cdot \widetilde{\tau}_{\text{resp}},
\quad
w_n = w_l = w_\tau = \tfrac{1}{3},
\label{eq:meth:lcog}
\end{equation}

\noindent где $\widetilde{N}_{\text{int}}$, $\widetilde{L}_{\text{ctx}}$
и $\widetilde{\tau}_{\text{resp}}$ --- нормированные на отрезок
$[0,1]$ значения числа обращений к участнику, средней длины
предоставляемого ему контекста и средней задержки ответа симулятора
соответственно. Нормировка выполняется по квантилям распределения
наблюдаемых значений за все запуски эксперимента. Равные веса
выбраны априорно до начала запусков как нейтральная свертка трех
независимых компонентов нагрузки.

Адаптивная конфигурация считается выигрышной относительно
наилучшей статической, если она дает прирост качества не менее
$5\%$ либо снижение стоимости не менее $15\%$ при сохранении
качества. Пороги фиксируются предварительно, до проведения
основных запусков, что соответствует практике предварительной
регистрации гипотез в эмпирических исследованиях.

\subsection{Корпус задач и обоснование его выбора}
\label{sec:meth:corpus}

Эксперименты проводятся на стратифицированной выборке из четырех
открытых корпусов, покрывающих четыре качественно различных
класса задач --- программирование, математические рассуждения,
творческая генерация и анализ данных. Сравнение корпуса
с альтернативами приведено в таблице~\ref{tab:meth:corpora}.

\begin{table}[!htbp]
\caption{Сравнение рассмотренных корпусов задач}
\label{tab:meth:corpora}
\footnotesize
\renewcommand{\arraystretch}{1.25}
\begin{tabularx}{\textwidth}{%
   >{\raggedright\arraybackslash}p{2.6cm}%
   >{\raggedright\arraybackslash}p{2.4cm}%
   c%
   >{\raggedright\arraybackslash}p{1.7cm}%
   >{\raggedright\arraybackslash}X}
\toprule
\textbf{Корпус} & \textbf{Класс} & \textbf{Размер} &
\textbf{Метрика} & \textbf{Решение по~включению} \\
\midrule
HumanEval~\cite{humaneval2021}        & Программирование    & 164   & pass@1        & Включен; стандарт сравнения. \\
HumanEval+~\cite{humanevalplus2023}   & Программирование    & 164   & pass@1+       & Не~включен; превышает покрытие основной выборки. \\
LiveCodeBench~\cite{livecodebench2024} & Программирование   & 500+  & pass@1        & Резерв для~проверки на~загрязнение. \\
GSM8K~\cite{gsm8k2021}                & Мат. рассуждения    & 8500  & exact match   & Включен; стандарт оценки. \\
MATH~\cite{mathbench2021}             & Мат. рассуждения    & 12500 & exact match   & Не~включен; сложность выше возможностей рабочих моделей. \\
MMLU-Pro~\cite{mmlupro2024}           & Множ. рассуждения   & 12000 & accuracy      & Не~включен; формат ответа не~подходит для~дебатной топологии. \\
CommonGen~\cite{commongen2020}        & Творческая генерация & 35000 & ROUGE-L $+$ cov & Включен; покрывает свободная (без~однозначного ответа) генерацию с~ограничениями. \\
InfiAgent-DABench~\cite{dabench2024}  & Анализ данных       & 257   & numeric exact & Включен; единственный корпус с~замкнутым ответом по~аналитике данных. Качество чувствительный к~модели (см.~\ref{sec:exp:conf}). \\
HellaSwag~\cite{hellaswag2019}        & Здравый смысл       & 10000 & accuracy      & Не~включен; качество современных моделей близко к~потолку. \\
\bottomrule
\end{tabularx}
\end{table}

Выбор обусловлен совокупностью требований. Каждый класс задач
представлен ровно одним корпусом, чтобы избежать смешения
эффектов корпуса и класса. Каждый корпус допускает автоматическую
объективную оценку, исключающую субъективность оценщика
(юнит-тесты, точное совпадение числа, метрики покрытия концептов).
Размер задач в каждом корпусе позволяет уложиться в один запуск
рабочих моделей. От рассматриваемых для каждого класса альтернатив
выбранные корпуса отличаются распространенностью в литературе и
зрелостью методики оценки, что обеспечивает сопоставимость
результатов с предшествующими работами. Из каждого корпуса
производится стратифицированная выборка по 15 задач, итого
60 задач в одном базовом запуске. Каждая задача повторяется
с тремя различными зернами генерации (для контроля стохастичности
выборки токенов), что в комбинации с пятью статическими топологиями
дает $60 \times 3 \times 5 = 900$ запусков в эксперименте E1.

\subsection{Языковые модели и инфраструктура}
\label{sec:meth:models}

Распределение моделей по ролям изолирует переменную «модель» от
переменных «топология» и «роль человека». Все рабочие агенты
(планировщик, исследователь, исполнитель, критик, дебатер), а также
маршрутизаторы и симулятор человека используют единую модель
Cerebras gpt-oss-120b (~3000 токенов/с на Cerebras WSE). Иерархия
«работник слабее критика» реализуется различиями в подсказках и
параметре строгости рассуждения, а не разными моделями ---
фиксированная модельная мощность критична для чистоты сравнения
топологий.

Судья оценки качества вынесен в отдельную линию --- OpenAI
gpt-4.1-mini с самосогласованностью из трех вызовов и нулевой
температурой; модель из другого семейства весов исключает петлю
положительной обратной связи. Подтверждающие запуски выполняются
на Cerebras qwen-3-235b как представителе другого семейства весов.

Выбор моделей опирался на четыре критерия --- отказоустойчивый
программный интерфейс (от англ. Application Programming Interface,
API) с вызовом инструментов и структурированным выводом, низкая
удельная стоимость вызова, высокая пропускная способность и
доступность из РФ.

Инфраструктура (глава~\ref{sec:architecture}) --- LangGraph,
PostgreSQL~16, Apache Parquet; параллелизм --- 12 процессов с
асинхронной обработкой внутри каждого.

\subsection{Дизайн экспериментов --- воронка}
\label{sec:meth:funnel}

Прямой эксперимент по всему декартову произведению (6 топологий
$\times$ 5 ролей $\times$ 3 режима $\times$ 4 класса задач) дал бы
более 10\,000 ячеек. Вместо этого исследование организовано как
воронка из четырех экспериментов, каждый из которых сужает
пространство конфигураций по результатам предыдущего
(таблица~\ref{tab:meth:funnel}).

\begin{table}[!htbp]
\caption{Структура экспериментальной воронки}
\label{tab:meth:funnel}
\small
\renewcommand{\arraystretch}{1.25}
\begin{tabularx}{\textwidth}{%
   l%
   c%
   >{\raggedright\arraybackslash}p{3.0cm}%
   >{\raggedright\arraybackslash}X}
\toprule
\textbf{Эксп.} & \textbf{Вопрос} & \textbf{Запуски} &
\textbf{Цель} \\
\midrule
E1 & RQ1 & $5\cdot4\cdot15\cdot3=900$ & Парето-фронт пяти статических
   топологий на~четырех классах задач. \\
\addlinespace[3pt]
E2 & RQ3 & $\approx 2295$ (после фильтра top-3) & Влияние пяти ролей
   человека на~каждую топологию из~top-3~E1. \\
\addlinespace[3pt]
E3 & RQ2 & $\approx 640$ (237 адаптивных $+$ 406 статических) &
   Адаптивная мета-топология в~двух режимах маршрутизации
   (\code{rule}, \code{llm}) и~три статические базовые линии. \\
\addlinespace[3pt]
E4 & RQ4 & $3\cdot4\cdot15\cdot3=540$ & Адаптивная роль человека
   в~сочетании с~адаптивной топологией. \\
\addlinespace[3pt]
Conf. & --- & $120+142\approx 260$ & Кросс-семейная валидация
   выводов E3 и~E4 на~подтверждающей модели Qwen-3-235B. \\
\bottomrule
\end{tabularx}
\end{table}

Исследовательские вопросы формулируются следующим образом.

\textbf{RQ1.} Различаются ли пять статических топологий
по~количественным показателям качества, стоимости и~времени
для~разных классов задач?

\textbf{RQ2.} Раскладывается на~две взаимодополняющих части. (а)
Дает ли адаптивная мета-топология значимый прирост качества или~экономию
стоимости по~сравнению с~per-task best-of референсом и~с~отдельными
статическими топологиями? (б) Устойчив ли выбор оптимальной
стратегии маршрутизации (правиловой или~языковой) к~смене семейства весов
рабочего агента, или~он является решение, принимаемое для~каждого семейства моделей?

\textbf{RQ3.} Какая из~пяти ролей человека наиболее эффективна для
каждой топологии и~для~каждого класса задач?

\textbf{RQ4.} Превосходит ли совместная адаптация топологии и~роли
человека по~фазам решения фиксированную роль; сохраняется ли
направление эффекта при~смене семейства весов рабочего агента?

Статистическая обработка включает двухфакторный дисперсионный
анализ (two-way ANOVA, тип~III сумм квадратов для несбалансированного
дизайна) с~эффектом топологии и~эффектом роли в~качестве факторов
и~их~взаимодействием на~уровне значимости $\alpha = 0{,}05$.
При~нарушении допущений нормальности или равенства дисперсий
применяется непараметрический аналог по~критерию Краскела--Уоллиса.
Парные сравнения средних выполняются с~поправкой
Бенджамини--Хохберга~\cite{benjamini1995} на~множественное сравнение
при~уровне ложных открытий $\text{FDR} = 0{,}05$. Оценка практической
значимости различий проводится через коэффициент размера эффекта
Коэна~$d$~\cite{cohen1988}. Доверительные интервалы (от англ.
Confidence Interval, CI) рассчитываются по~методу непараметрического
бутстрапа~\cite{efron1979} (метод процентилей) с~десятью тысячами
ресэмплов на~уровне 95\,\% при~фиксированном начальном значении
генератора псевдослучайных чисел.

\subsection{Разграничение заимствованных и авторских элементов}
\label{sec:meth:boundary}

Граница между заимствованными и авторскими элементами. К
заимствованным относятся пять статических топологий, четыре
открытых корпуса задач, языковые модели и стандартный
статистический инструментарий (ANOVA, бутстрап, коэффициент Коэна).
К авторским --- адаптивная мета-топология вместе с процедурой
работы (\ref{eq:meth:adaptive}); формализация пяти ролей человека
и матрица их совместимости с топологиями; набор адаптивно-специфичных
метрик ($N_{\text{switch}}$, $r_{\text{guard}}$, $\gamma$,
$r_{\text{hurt}}$, $\Delta_{\text{oracle}}$); вороночная структура
исследования с предварительной фиксацией порогов
(\ref{eq:meth:thresholds}) и кросс-семейным подтверждением.

\subsection{Выводы и результаты по главе}
\label{sec:meth:summary}

Формализована MAS как кортеж из шести компонентов с динамическим
переключением графа по фазе; определены шесть топологий, пять
ролей человека и три группы метрик (включая адаптивно-специфичные).
Обоснован выбор корпуса из четырех бенчмарков и линии моделей с
кросс-семейным судьей. Сформулированы четыре RQ и воронка из
четырех экспериментов с подтверждающим запуском на дополнительном
семействе моделей. Граница заимствованного и авторского
зафиксирована в~\ref{sec:meth:boundary} и развернута
в~\ref{sec:contribution}. План реализуется в инфраструктуре
главы~\ref{sec:architecture} и проверяется в~\ref{sec:experiments}.
